\documentclass[11pt]{article}

\usepackage[margin=1in]{geometry}
\usepackage{amsmath,amssymb,amsthm}
\usepackage{xcolor}
\usepackage[hidelinks]{hyperref}

\newtheorem{theorem}{Theorem}[section]
\newtheorem{lemma}[theorem]{Lemma}
\newtheorem{corollary}[theorem]{Corollary}
\theoremstyle{definition}
\newtheorem{definition}[theorem]{Definition}
\theoremstyle{remark}
\newtheorem{remark}[theorem]{Remark}

\newcommand{\cut}{\delta}

\title{Four universally separated marks in\\
Tait-colourable cubic graphs}
\author{AI-assisted research draft for independent human verification\\
\small Human author attribution to be supplied before submission}
\date{26 July 2026}

\begin{document}
\maketitle

\noindent
\colorbox{yellow!22}{%
  \parbox{\dimexpr\linewidth-2\fboxsep\relax}{%
  \textbf{AI-assistance disclosure.}
  OpenAI Codex agents (GPT-5-series models, accessed 26 July 2026)
  formulated the marked theorem, found the proof architecture, supplied
  the cut and decomposition arguments, checked the cited source
  statements, and drafted this manuscript under human direction.
  An additional Codex agent performed the publication audit, but that is
  not independent human verification or peer review.  No human author or
  referee is represented as having certified the theorem, proof, or
  novelty search.  This is a research draft, not a submission-ready
  paper.  A human author must check every reduction and source, conduct
  an independent literature review, and accept ordinary scholarly
  responsibility before public submission.}}

\begin{abstract}
Let \(H\) be a connected simple cubic graph admitting a proper
three-edge-colouring, and let \(S\) be a matching of four edges.
Suppose that in every proper three-edge-colouring, no bichromatic cycle
contains two members of \(S\).  Suppose also that every vertex set
\(X\) for which \(H[X]\) contains a cycle satisfies
\[
 |\cut_H(X)|+|S\cap E(H[X])|\ge4.
\]
We prove that \(H\) has an even edge-subgraph containing \(S\) such that
every cycle component contains an even number of members of \(S\).
More precisely, the proof produces either one cycle through all four
marks or two disjoint cycles containing two marks each.

The proof first resolves cyclic cuts of size two and all marked cyclic
cuts of size three.  In the remaining case, a self-contained cubic
three-cut decomposition yields cyclically \(4\)-edge-connected atoms.
A central atom is selected by the distribution of the four marks.
Kempe switching gives a Tait colouring in which all marks have one
colour.  The prescribed-edge theorems of Aldred, Ellingham, Hemminger,
and Holton and of Knappe and Pitz supply the central cycle, according
to whether the four selected edges remain distinct.  One-mark paths
lift it through the branches.

The argument is computer-free.  It proves the stated marked-graph
theorem only; it neither proves nor refutes the five-cycle double cover
conjecture.  The exact novelty claim is provisional pending independent
expert review.
\end{abstract}

\section{Introduction}

A \emph{Tait colouring} of a cubic graph is a proper edge-colouring with
three colours.  For a fixed Tait colouring, each pair of colours induces
a disjoint union of cycles.  We call a matching \(S\)
\emph{universally separated} when, in every Tait colouring, each such
bichromatic cycle contains at most one edge of \(S\).

The theorem below arose in a reduction related to five-cycle double
covers.  That motivation is not used in the proof.  In particular, the
standard five-cycle double cover conjecture remains a separate open
problem; see Oum's formulation as Conjecture~18
\cite[Section~9.1]{oum2026}.

\begin{theorem}[four-mark theorem]\label{thm:main}
Let \(H\) be a connected simple Tait-colourable cubic graph, and let
\(S\subseteq E(H)\) be a universally separated matching of size four.
Assume that
\begin{equation}
 |\cut_H(X)|+|S\cap E(H[X])|\ge4                 \tag{M}
\end{equation}
whenever \(H[X]\) contains a cycle.  Then \(H\) has a binary cycle
\(Q\subseteq E(H)\) such that
\begin{enumerate}
  \item \(S\subseteq Q\); and
  \item every component of \(Q\) contains an even number of edges of
        \(S\).
\end{enumerate}
More precisely, \(Q\) can be chosen either as one cycle containing all
four marks or as the union of two disjoint cycles, each containing two
marks.
\end{theorem}

Here a \emph{binary cycle} means an edge-subset of even degree at every
vertex.  In a cubic graph its nonempty components are cycles.

The two external cycle results used below are exact and classical.
Knappe and Pitz prove, in particular, that in a \(3\)-edge-connected
graph a set of at most three edges lies in a circuit if and only if it
contains no odd edge cut \cite[Fact~5.4]{knappe2020}.  Aldred,
Ellingham, Hemminger, and Holton prove that four independent edges of a
quasi \(4\)-connected graph lie on a common cycle
\cite[Theorem~3.3]{aldred1997}.  They also record that cyclically
\(4\)-edge-connected cubic graphs are quasi \(4\)-connected.

Related general tree decompositions of \(3\)-connected graphs into
cyclically \(4\)-edge-connected components were developed by Pastor
\cite{pastor2018}; the exact cubic construction also appears in
Nedela, Seifrtov\'a, and \v Skoviera \cite[Section~6]{nedela2024}.
To avoid any convention or hypothesis gap, we prove the elementary
cubic specialization needed here.

We found no prior statement of Theorem~\ref{thm:main} in the sources
above or in targeted searches using its marked-colouring and cut
hypotheses.  This is only a provisional novelty assessment.  The
prescribed-edge cycle theorems, Kempe switching, cut parity, and
three-cut decomposition are not claimed as new.

\section{Colour flexibility and pole lemmas}

\begin{lemma}[mark precolouring]\label{lem:precolour}
Let \(S\) be universally separated in a Tait-colourable cubic graph.
Every map from \(S\) to the three colours extends to a Tait colouring.
\end{lemma}

\begin{proof}
Start with any Tait colouring.  Suppose that a mark \(e\) currently has
colour \(p\) and is prescribed colour \(q\ne p\).  The unique
\(pq\)-bichromatic cycle through \(e\) contains no other mark.
Interchanging \(p\) and \(q\) on that cycle changes the colour of \(e\)
without changing another member of \(S\).  Process the marks one at a
time.  Universal separation applies after every switch, so a later
switch does not disturb a mark already processed.
\end{proof}

\begin{lemma}[cut and cap facts]\label{lem:capfacts}
Let \(G\) be a simple Tait-colourable cubic graph.
\begin{enumerate}
  \item \(G\) is bridgeless.
  \item In a Tait colouring, the three edges of a three-edge cut have
        three distinct colours.
  \item If \(G\) has no edge cut of size at most two and \(B\) is a
        cyclic three-edge cut, then the three edges of \(B\) have
        distinct ends on each shore.
  \item Capping either shore of \(B\) by one new vertex adjacent to its
        three boundary ends produces a simple cubic
        \(3\)-edge-connected graph.
\end{enumerate}
\end{lemma}

\begin{proof}
Every edge lies on either of the two bichromatic cycles using its colour,
which proves (1).  For (2), if \(X\) is a shore, then the number of
edges of each colour in \(\cut_G(X)\) has the same parity as \(|X|\).
A cubic shore of a three-edge cut has odd order, so each colour occurs
oddly; all three therefore occur once.

For (3), suppose two cut edges have a common end \(v\) on a shore \(A\).
The sole edge from \(v\) into \(A-v\), together with the third cut edge,
separates \(A-v\).  Every cycle of \(G[A]\) avoids \(v\), since \(v\)
has only one neighbour in \(A\).  This gives a cyclic cut of size at
most two, a contradiction.

The cap is simple by (3).  If a set \(Y\) defines a cut of size at most
two in the cap, choose the shore not containing the new cap vertex.
Replacing cap edges by the corresponding members of \(B\) gives an
edge cut of the same size in \(G\).  This proves (4).
\end{proof}

We use a \emph{\(k\)-pole} only as shorthand for a connected induced
shore of a \(k\)-edge cut, with the cut edges regarded as boundary
semiedges.  The boundary ends below are distinct.

\begin{lemma}[two-mark shore]\label{lem:twomark}
Let \(P=H[A]\) be a connected two- or three-pole of minimum degree two.
Suppose exactly two marks lie internally in \(P\), and (M) holds for
every cyclic vertex set contained in \(A\).  Then one cycle of \(P\)
contains both marks.
\end{lemma}

\begin{proof}
Delete the bridges of \(P\), and form the tree of bridge-components.
If \(P\) has no bridge, then it is \(2\)-edge-connected.  A
\(2\)-edge-connected subcubic graph of minimum degree two has no
cutvertex: two components behind a cutvertex would each require at
least two incident edges at that vertex.  Thus \(P\) is
\(2\)-connected.  Subdivide the two marks.  Two vertices of a
\(2\)-connected graph lie on a common cycle, which gives the required
cycle after undoing the subdivisions.

Now suppose the bridge-component tree is nontrivial.  Every leaf
component \(K\) is cyclic.  It cannot be an isolated vertex, since it
has one incident bridge but every vertex of \(P\) has degree at least
two.  Let \(p(K)\) count pole boundary edges incident with \(K\), and
let \(s(K)\) count its internal marks.  Exactly one bridge leaves \(K\),
so (M) gives
\[
 1+p(K)+s(K)\ge4.
\]
If \(K\) does not contain both marks, then \(s(K)\le1\), and hence
\(p(K)\ge2\).  Two such leaves would consume at least four boundary
edges, impossible for a two- or three-pole.  Some leaf component
therefore contains both marks.  It is \(2\)-connected by the preceding
argument and has a cycle through them.
\end{proof}

\begin{lemma}[one-mark shore]\label{lem:onemark}
Let \(P=H[A]\) be a connected three-pole of minimum degree two with
boundary ends \(u_1,u_2,u_3\).  Suppose exactly one mark \(e\) lies
internally in \(P\), and (M) holds for cyclic vertex sets contained in
\(A\).  For every pair \(i\ne j\), \(P\) has a
\(u_i\)-to-\(u_j\) path containing \(e\).
\end{lemma}

\begin{proof}
The pole has no bridge.  Otherwise, a leaf \(K\) of its
bridge-component tree is cyclic as in Lemma~\ref{lem:twomark}.  The
inequality
\[
 1+p(K)+s(K)\ge4
\]
says that an unmarked leaf consumes all three boundary edges and a
marked leaf consumes at least two.  Any two leaves therefore consume
more than the three available boundary edges.

Thus \(P\) is \(2\)-edge-connected and hence \(2\)-connected by the
subcubic cutvertex argument.  Add a new edge \(u_i u_j\), allowing a
parallel edge temporarily if necessary.  In a \(2\)-connected graph,
any two edges lie on a common cycle.  A cycle through the new edge and
\(e\), with the new edge removed, is the required path.
\end{proof}

\section{Low cyclic cuts}

\begin{lemma}[cyclic two-cuts]\label{lem:twocut}
If \(H\) satisfies the hypotheses of Theorem~\ref{thm:main} and has a
cyclic two-edge cut, then the conclusion of the theorem holds.
\end{lemma}

\begin{proof}
Apply (M) to both shores.  Each shore contains at least two internal
marks.  Since there are only four marks, the cut is unmarked and each
shore contains exactly two.  Lemma~\ref{lem:twomark} gives a cycle
through the two marks on each shore.  Their disjoint union is the
required binary cycle.
\end{proof}

\begin{lemma}[cyclic three-cuts]\label{lem:threecut}
Let \(H\) satisfy the hypotheses of Theorem~\ref{thm:main}, have no
cyclic two-edge cut, and let \(B\) be a cyclic three-edge cut.
Unless \(B\) is unmarked and splits the internal marks \(1+3\), the
conclusion of Theorem~\ref{thm:main} holds.
\end{lemma}

\begin{proof}
Let \(r=|B\cap S|\), and let \(s_A,s_D\) count internal marks on the two
shores.  Inequality (M) gives \(s_A,s_D\ge1\), while
\[
 r+s_A+s_D=4.
\]
Hence \(r\le2\).

If \(r=2\), each shore contains one internal mark.  The two marked
members of \(B\) determine the same boundary pair on both shores.
Lemma~\ref{lem:onemark} supplies a path through the internal mark on
each side for that pair.  The paths and the two marked cut edges glue
to a cycle through all four marks.

Suppose \(r=1\), with two internal marks on \(A\) and one on \(D\).
Cap \(A\) by a new vertex \(z\).  The cap is simple and
\(3\)-edge-connected by Lemma~\ref{lem:capfacts}.  Let \(F\) consist of
the cap edge corresponding to the marked member of \(B\), together
with the two internal marks of \(A\).  These are three independent
edges.

The set \(F\) contains no odd edge cut of the cap.  Indeed, edge
connectivity leaves only the possibility that \(F\) itself is a
three-edge cut.  Since \(F\) is independent, this cut is nontrivial and
both shores are cyclic: a connected shore of order \(n>1\) has
\((n-1)/2\ge1\) independent cycles by the cubic degree sum.  Choose its
shore \(X\) not containing \(z\).
Viewed in \(H\), it has \(|\cut_H(X)|=3\), while both internal marks of
\(A\) cross the cut and no mark lies internally in \(X\).  This
contradicts (M).

Knappe and Pitz's \(g(3)=3\) result
\cite[Fact~5.4]{knappe2020} now gives a circuit through \(F\).
Because the cap is cubic, the circuit is a cycle.  At \(z\) it uses the
marked cap edge and one other cap edge.  Deleting \(z\) leaves a path
through the two internal marks for that boundary pair.
Lemma~\ref{lem:onemark} supplies a path through the sole internal mark
of \(D\) for the same pair.  Gluing the paths through the corresponding
two members of \(B\) gives a cycle containing all four marks.

It remains that \(r=0\).  The internal split is \(1+3\) or \(2+2\).  In
the \(2+2\) case, Lemma~\ref{lem:twomark} on both shores gives two
closed cycles whose union proves the theorem.  The only unresolved
local case is therefore an unmarked \(1+3\) cut.
\end{proof}

\section{A cubic atom tree}

For a cyclic three-edge cut, \emph{capping a shore} means adjoining one
new cubic vertex at its three boundary ends.  Conversely, two capped
graphs with specified cap vertices can be joined by deleting the cap
vertices and matching their three incident ports.  We call this a
\emph{three-sum}.  The matching of ports is retained as part of the
data.

\begin{lemma}[cubic atom tree]\label{lem:atoms}
Let \(G\) be a simple cubic graph with no edge cut of size at most two.
There is a finite tree \(T\) with a simple cubic graph \(G_t\) at every
node \(t\) such that:
\begin{enumerate}
  \item each \(G_t\) is \(3\)-edge-connected and has no cyclic
        three-edge cut;
  \item \(G\) is recovered by iterated three-sums along the edges of
        \(T\); and
  \item each edge of \(T\) corresponds, after expanding all other
        factors, to a cyclic three-edge cut of \(G\).
\end{enumerate}
Every \(G_t\) is a cyclically \(4\)-edge-connected cubic graph and hence
is quasi \(4\)-connected.
\end{lemma}

\begin{proof}
Proceed recursively.  If \(G\) has no cyclic three-edge cut, use the
one-node tree.  Otherwise choose such a cut, cap its two shores, and
recurse on the two capped graphs.  Lemma~\ref{lem:capfacts} shows that
both are simple, cubic, and \(3\)-edge-connected.  Each has fewer
vertices than \(G\), so the recursion terminates.

In each recursive tree, the cap vertex introduced by the chosen cut
belongs to one terminal factor.  Join those two factor nodes by a new
tree edge and remember the port matching.  Reversing this operation
recovers \(G\), proving (2).

For (3), expand the caps on either side of a recursive cut.  Put an
expanded branch on the same side as its cap vertex.  Every cap edge in
the cut is replaced by exactly one original edge, and all other port
edges remain internal.  Thus the boundary still has size three.
A cycle using a cap vertex lifts through a path in the connected
expanded branch between the two used ports; a cycle avoiding the cap
is unchanged.  Cyclicity of both shores is therefore preserved.

It remains only to justify the last sentence.  A simple
\(3\)-edge-connected cubic graph with no cyclic three-edge cut is
\(3\)-vertex-connected.  A cutvertex would force two components each
to use at least two of its three incident edges.  If \(\{u,v\}\) were
a two-vertex cut, then \(u\) and \(v\) could not be adjacent; the six
edges from them would split into two three-edge boundaries.  Simplicity
forces both corresponding shores to contain cycles, contrary to the
absence of a cyclic three-edge cut.  The factor is consequently
cyclically \(4\)-edge-connected in the standard cubic sense.  Aldred
et al.\ observe that every such cubic graph is quasi
\(4\)-connected \cite{aldred1997}.
\end{proof}

\begin{lemma}[central node]\label{lem:central}
Let nonnegative integer weights of total four be assigned to the
vertices of a finite tree \(T\).  Suppose every tree edge separates the
total weight as \(1+3\).  Then \(T\) has a vertex \(a\) such that every
component of \(T-a\) has weight one.  If \(a\) has weight \(r\) and
degree \(q\), then \(r+q=4\).
\end{lemma}

\begin{proof}
Orient every edge toward its side of weight three.  A finite directed
tree has a sink \(a\).  The side away from \(a\) across each incident
edge has weight one.  These are exactly the components of \(T-a\), so
their number is \(q\) and \(r+q=4\).
\end{proof}

\section{Proof of the four-mark theorem}

\begin{proof}[Proof of Theorem~\ref{thm:main}]
Suppose for a contradiction that the required binary cycle does not
exist.  Lemma~\ref{lem:twocut} eliminates cyclic two-edge cuts.  A
Tait-colourable cubic graph is bridgeless, and every two-edge cut in a
simple bridgeless cubic graph is cyclic.  Indeed, a minimal two-edge cut
has connected shores, and a shore of order \(n\) has
\[
 |E(H[X])|-|X|+1=n/2>0
\]
by the cubic degree sum; the same holds on the other shore.  Hence
\(H\) has no edge cut of size at most two.

Apply Lemma~\ref{lem:atoms}.  Every tree edge expands to a cyclic
three-edge cut of \(H\).  By Lemma~\ref{lem:threecut}, each such cut is
unmarked and splits the four marks \(1+3\).  Since no decomposing cut is
marked, every mark belongs internally to a unique atom.  Weight each
atom by its number of marks.  Lemma~\ref{lem:central} gives a central
atom \(A\), containing \(r\) internal marks and incident with
\(q=4-r\) branches, each containing exactly one mark.

By Lemma~\ref{lem:precolour}, choose a Tait colouring of \(H\) in which
all four marks have one colour, say \(c\).  The three edges of every
decomposing cut have distinct colours.  The colouring therefore induces
a Tait colouring on every capped atom by retaining those colours on the
three cap edges.

In \(A\), select its \(r\) internal marked edges and, at each of its
\(q\) cap vertices, select the unique incident edge of colour \(c\).
Let \(F\) be the set of distinct selected edges.  Because all members of
\(F\) have colour \(c\), \(F\) is a matching; it contains every internal
mark, meets every cap vertex, and has size at most four.
Two cap vertices can be adjacent.  In that case their unique incident
\(c\)-edge may be the same edge, which is why \(F\) is defined as a set
of distinct edges and may have size less than four.  A cycle containing
that shared edge nevertheless visits both cap vertices.

There is a cycle \(C\) of \(A\) containing \(F\).  If \(|F|=4\), this
is the independent-edge theorem of Aldred, Ellingham, Hemminger, and
Holton for quasi \(4\)-connected graphs
\cite[Theorem~3.3]{aldred1997}.  If
\(|F|\le3\), then \(F\) contains no odd cut: a one-edge cut is excluded
by \(3\)-edge-connectivity, and a three-edge cut contained in a
matching would be nontrivial and cyclic, which does not occur in an
atom.  Knappe and Pitz's \(g(3)=3\) theorem
\cite[Fact~5.4]{knappe2020} gives a circuit through \(F\), and in
the cubic graph \(A\) its edge set is a cycle.

The cycle \(C\) visits every cap vertex, since it contains a selected
edge incident with that vertex.  At a cap vertex it uses exactly two
ports.  Expand the corresponding component of the atom tree.  In the
original graph this component is a connected three-pole containing
exactly one mark.  It satisfies (M) on every cyclic subset, so
Lemma~\ref{lem:onemark} gives a path through its mark between the two
ports used by \(C\).  Replace the two-edge passage through the cap
vertex by that path and the corresponding boundary edges.

Perform this replacement for all \(q\) branches.  Formally, first
reconstruct each component of \(T-A\) internally while retaining its
root cap, and then attach those branch composites to \(A\) one at a
time.  If two cap vertices of \(A\) are adjacent, expanding the
first deletes their shared cycle edge and reconnects its dangling end
to the inserted branch.  The resulting edge is still the same port of
the lifted cycle at the second cap vertex; that vertex remains cubic and
on the cycle.  Expanding the second cap later replaces its two current
ports in exactly the same way.  Thus adjacency of cap vertices causes
no ambiguity and no splitting of the cycle.

Each sequential replacement of a cap passage by the corresponding
one-mark path preserves a single cycle.  The final cycle lies in \(H\)
and contains the \(r\) central marks and the one mark from each of the
\(q\) branches, hence all four members of \(S\).  This contradicts the
assumed failure and proves the theorem.
\end{proof}

\begin{remark}[scope of the one-cycle conclusion]\label{rem:onecycle}
The central-atom argument and the marked three-cut cases construct one
cycle through all four marks.  The proof may instead terminate at an
unmarked cyclic three-cut with a \(2+2\) mark split, where
Lemma~\ref{lem:threecut} produces a binary cycle with two components.
Thus this manuscript does \emph{not} claim the stronger statement that
one cycle always contains all four marks.  The precise alternative
proved is one four-mark cycle or two disjoint two-mark cycles.
\end{remark}

\section*{AI-use statement}

The page-one disclosure is part of this draft and should remain in every
derivative version.  Codex agents materially contributed to the theorem
statement, proof search, low-cut lemmas, atom-tree proof, source audit,
and exposition.  The proof is written for direct human checking and
uses no finite computation, SAT solver, or unverifiable certificate.
That makes it human-checkable; it does not mean a human has checked it.
No independent human verification is asserted.

The precise source-dependent inputs are the prescribed-circuit result
of Knappe and Pitz \cite{knappe2020}, the quasi-\(4\)-connected
prescribed-cycle result of Aldred, Ellingham, Hemminger, and Holton
\cite{aldred1997}, and, for context only, the broader decomposition
theory \cite{pastor2018,nedela2024}.  Any prospective human author should verify
those applications against the published papers and repeat the novelty
search before submission.

\bibliographystyle{plain}
\bibliography{references}

\end{document}
